\chapter{Conclusions and Future Work}
\label{ch:con}
\section{Conclusions}
This project aimed to answer a focused but societally important question: Can year-over-year Zillow Home Value Index (ZHVI) growth be predicted accurately from historical price patterns and ZIP-code-level demographics? In tackling that question, we first clarified why it matters. Reliable ZHVI forecasts give homebuyers, investors, planners, and policy-makers an early signal of local economic health and help guide billions of dollars in housing decisions. Yet the task is notoriously hard because housing markets are “noisy,” highly non-linear, and riddled with missing or repeated observations across geography and time. To meet these challenges, we engineered a pipeline that begins with rigorous data curation, combining Zillow price histories with ACS demographic variables, pruning long NA streaks, and one-hot encoding location. We then enforced a ZIP-based train/test split to avoid data leakage, standardized the features. Three supervised models: Decision Tree, K-Nearest Neighbors, and Random Forest, were tuned significantly. The results show a clear winner: the tuned Random Forest delivered the lowest test-set error and the most stable rankings across bootstrap runs. Equally important, the feature-selection step improved every model, proving that careful variable curation is as valuable as algorithm choice when signal-to-noise is low. The exercise also surfaced key domain insights. ZIP-level volatility remained high even after cleaning, confirming the heterogeneous nature of local housing markets and underscoring why global heuristics so often fail. Beyond the metrics, the project’s main contributions are methodological. We provide a reproducible template for integrating public census data with proprietary price indices, demonstrate a leakage-proof geographic split strategy, and illustrate how mutual-information scores can stabilize feature importance in real-estate time series. Practically, stakeholders now have an interpretable model that highlights which demographic and historical factors precede outsized appreciation—knowledge that can inform targeted investment, zoning debates, and equity-oriented housing policy. The work is not without limits. Predictions are confined to a one-year horizon and ZIP codes with complete historical records; extrapolating farther in time or space will require additional exogenous signals and more granular economic covariates. Future iterations should therefore one, extend the framework to multi-year forecasts and two, stream in higher-frequency external data—mortgage rates, permits, or employment shocks—to sharpen responsiveness to market inflections. In sum, the project validates that with thoughtful preprocessing, robust feature selection, and ensemble learning, it is possible to cut through the noise of local housing data and deliver actionable, ZIP-specific ZHVI growth estimates.

\section{Future work}
This section should refer to Chapter~\ref{ch:results} where the author has reflected their criticality about their own solution. Concepts for future work are then sensibly proposed in this section.

\textbf{Guidance on writing future work:} While working on a project, you gain experience and learn the potential of your project and its future works. Discuss the future work of the project in technical terms. This has to be based on what has not been yet achieved in comparison to what you had initially planned and what you have learned from the project. Describe to a reader what future work(s) can be started from the things you have completed. This includes identifying what has not been achieved and what could be achieved. 



A good future work summary could be approximately 300--500 words long, but this is just a recommendation.
